% SYNTHETIC TEST UNIT (integration tools) -- deliberate errors are marked "DELIBERATE"
\chapter{减除法（合成测试）}\label{app:appN}

\section{假设}\label{sec:appN-asm}

% DELIBERATE C-P10 W (the 被使用于 field lists only one of the users)
\begin{assumption}[Collins 条件 \tagA]\label{asm:appN-collins}
(CS1)--(CS3)（合成测试）。
地位：全阶论证覆盖。
文献：Collins 2011。
被使用于：\cref{thm:appN-subtraction}。
解除途径：一般图的证明。
\uses{def:dis-kinematics}
\end{assumption}

\begin{assumption}[取值规则]\label{asm:appN-values}
(V1)(V2)（合成测试）。
地位：约定。
文献：无。
被使用于：\cref{thm:appN-subtraction,thm:method-template}。
解除途径：无（约定）。
% uses: none
\end{assumption}

\begin{assumption}[BRST 分类 \tagC]\label{asm:appN-brst-nonlocal}
合成测试。
地位：无全阶论证（具名缺口）。
文献：Collins 2011。
被使用于：\cref{thm:appN-structure,prop:dy-ctech}。
解除途径：分类定理。
% uses: none
\end{assumption}

% DELIBERATE C-P9 E assumption used by no statement
\begin{assumption}[未被使用的假设 \tagA]\label{asm:appN-unused}
合成测试。
地位：全阶论证覆盖。
文献：无。
被使用于：无。
解除途径：无。
% uses: none
\end{assumption}

% DELIBERATE C-P2/K5 E register [P], written [A] (strengthened); K14 W fields 文献/解除途径 missing
\begin{assumption}[标签不符的假设 \tagA]\label{asm:appN-wrongtag}
合成测试。
地位：仅在特例中证明。
被使用于：\cref{thm:appN-structure}。
% uses: none
\end{assumption}

\section{定理}\label{sec:appN-thm}

\begin{theorem}[条件减除定理 \tagR]\label{thm:appN-subtraction}
（$R_{\mathrm{cond}}$：在所列编号假设之下证明完整；结论的地位不高于其中最弱的假设。）前提：\cref{asm:appN-collins}、\cref{asm:appN-values}。
合成测试。（本定理只涉及部分子态。）
\uses{asm:appN-collins,asm:appN-values,def:dis-kinematics}
\end{theorem}
\begin{proof}略。\end{proof}

% DELIBERATE C-P8 E the premise sentence misses asm:appN-wrongtag (in \uses) and lists lem:appN-cyc-a (not in \uses)
\begin{theorem}[因子化结构 \tagR]\label{thm:appN-structure}
（$R_{\mathrm{cond}}$：在所列编号假设之下证明完整；结论的地位不高于其中最弱的假设。）前提：\cref{asm:appN-brst-nonlocal}、\cref{thm:appN-subtraction}（及 \cref{lem:appN-cyc-a}）。
合成测试。（本定理只涉及部分子态。）
\uses{asm:appN-brst-nonlocal,asm:appN-wrongtag,thm:appN-subtraction}
\end{theorem}
\begin{proof}略。\end{proof}

% DELIBERATE C-P8/K7 E no premise sentence
\begin{theorem}[缺前提句 \tagR]\label{thm:appN-noprem}
（$R_{\mathrm{cond}}$：在所列编号假设之下证明完整；结论的地位不高于其中最弱的假设。）
合成测试。（本定理只涉及部分子态。）
\uses{asm:appN-collins}
\end{theorem}

% DELIBERATE C-P8/K7 E the premise sentence is not among the first two sentences
\begin{proposition}[前提句靠后 \tagR]\label{prop:appN-late}
（$R_{\mathrm{cond}}$：在所列编号假设之下证明完整；结论的地位不高于其中最弱的假设。）这是第二句。前提：\cref{asm:appN-collins}。
合成测试。（本命题只涉及部分子态。）
\uses{asm:appN-collins}
\end{proposition}

% DELIBERATE C-P4/K4 E dependency cycle (and C-P3 E: uses a later statement of the same unit)
\begin{lemma}[环 A]\label{lem:appN-cyc-a}
合成测试。
\uses{lem:appN-cyc-b}
\end{lemma}

\begin{lemma}[环 B]\label{lem:appN-cyc-b}
合成测试。
\uses{lem:appN-cyc-a}
\end{lemma}

% (used by lem:dis-hidden of Ch07: hidden forward dependency on Ch12)
\begin{lemma}[附录中的辅助引理]\label{lem:appN-hidden}
合成测试。
\uses{thm:dy-factorization}
\end{lemma}
